\section{Experiments}
\label{sec:experiments}

\subsection{Setup}

We test \ours{} on 9 benchmarks spanning three modalities (EEG, images,
text) and multiple architectures:

\begin{itemize}[leftmargin=*]
\item \textbf{FACED} \cite{liu2024faced}: 9-class EEG emotion classification,
10{,}332 samples (6{,}720 train / 1{,}680 val / 1{,}932 test).
Architecture: EMOD axial-transformer with depth=6 (3 pretrained +
3 random init), $n_\mathrm{filters}=128$.
\item \textbf{SEED-V} \cite{zheng2019investigating}: 5-class EEG emotion
classification. Same EMOD d6 architecture.
\item \textbf{CIFAR-10}, \textbf{CIFAR-100} \cite{krizhevsky2009learning}:
image classification with a small CNN (64$\to$128$\to$128$\to$256 channels,
$\sim$500K params).
\item \textbf{MNIST}, \textbf{Fashion-MNIST}: with small MLP (3 layers,
FC-FC-FC) and small CNN respectively.
\item \textbf{SVHN}: house number recognition with small CNN.
\end{itemize}

For each dataset, we train 5 seeds with $T_A$ epochs and 5 seeds with
$T_B$ epochs (same augmentation, same optimizer, same LR schedule except
for $T_\mathrm{max}$), then compare:
(a)~individual-seed accuracy,
(b)~5-seed same-length deep ensemble accuracy (gold-standard baseline),
(c)~10-model mixed-length \ours{} ensemble accuracy.

\subsection{EEG Emotion Classification}

\begin{table}[h]
\centering
\caption{FACED 9-class EEG emotion classification. EMOD d6 axial transformer,
aug+LS+rand9 KD recipe. Our 10-model mixed-length ensemble is a new SOTA,
$+5.3$~pp over the EMOD AAAI 2026 baseline.}
\label{tab:faced}
\begin{tabular}{lcc}
\toprule
Method & Balanced Acc & $\Delta$ over baseline \\
\midrule
REVE NeurIPS 2025 \cite{kostas2025reve} & 0.500 & $-$0.142 \\
EMOD AAAI 2026 \cite{chen2026emod} & 0.642 & --- \\
EMOD d3 + aug + LS=0.1 + rand9 KD (ours, cycle 72) & 0.6467 $\pm$ 0.007 & $+$0.005 \\
EMOD d6 (3 pretrained + 3 random) single model & 0.6581 $\pm$ 0.007 & $+$0.016 \\
EMOD d6 5-seed ensemble & 0.6798 & $+$0.038 \\
EMOD d6 10-seed same-length ensemble & 0.6838 & $+$0.042 \\
\rowcolor{blond}
\textbf{EMOD d6 + d6\_e150 10-model \ours{}} & \textbf{\facedacc{}} & \textbf{$+$0.053} \\
\bottomrule
\end{tabular}
\end{table}

\begin{table}[h]
\centering
\caption{SEED-V 5-class EEG emotion classification. EMOD d6 cross-dataset
transfer. No 150ep checkpoints available, so we only report the standard
5-seed ensemble gain.}
\label{tab:seedv}
\begin{tabular}{lc}
\toprule
Method & Balanced Acc \\
\midrule
EMOD d3 + aug + KD (ours, cycle 244) & 0.3744 \\
EMOD d6 + aug + rand9 KD, individual mean & 0.3743 $\pm$ 0.006 \\
\rowcolor{blond}
\textbf{EMOD d6 5-seed ensemble} & \textbf{\seedvacc{}} (\seedvgain{}) \\
\bottomrule
\end{tabular}
\end{table}

The EEG emotion SOTA is a concrete flagship application: our ensemble trick
pushes FACED from $0.642$ to \facedacc{}, an $+8.2\%$ relative improvement.

\subsection{Image Classification (Small CNN and MLP)}

We next test whether the ensemble trick generalizes beyond the EEG
domain. We use standard small-CNN and MLP architectures on five vision
benchmarks covering a range of individual-model accuracies from below
$0.70$ (CIFAR-100) to above $0.98$ (MNIST).

\begin{table}[h]
\centering
\caption{\ours{} on four image datasets with small CNN/MLP architectures.
Each row: 5 seeds × 50ep (group A) + 5 seeds × 75ep (group B), mixed 10-model
ensemble. Gain scales inversely with individual accuracy.}
\label{tab:image}
\begin{tabular}{lcccc}
\toprule
Dataset & Arch & Ind.\ mean & 5-seed ensemble & \textbf{10-model mixed} \\
\midrule
CIFAR-10 & Small CNN & 0.9015 & 0.9209 (+0.0194) & \textbf{\ciftenacc{} (+0.0238)} \\
CIFAR-10 & ResNet-18 & 0.9452 & 0.9563 (+0.0111) & \textbf{0.9599 (+0.0147)} \\
CIFAR-100 & Small CNN & 0.6886 & 0.7225 (+0.0339) & \textbf{0.7353 (+0.0467)} \\
CIFAR-100 & ResNet-18 & 0.7707 & 0.8020 (+0.0313) & \textbf{0.8104 (+0.0397)} \\
MNIST & MLP & 0.9849 & 0.9859 (+0.0010) & \textbf{0.9863 (+0.0014)} \\
Fashion-MNIST & Small CNN & 0.9313 & 0.9418 (+0.0105) & \textbf{0.9434 (+0.0121)} \\
SST-2 & DistilBERT & 0.9062 & 0.9094 (+0.0032) & 0.9094 (+0.0032) \\
SVHN$^{\dagger}$ & Small CNN & -- & -- & -- \\
\bottomrule
\end{tabular}
\\[0.4em]
{\footnotesize $\dagger$: SVHN run excluded from main table. The short-CNN
trainer used for this sweep had a training instability that caused 7 of
10 seeds to fail to escape the uniform-prediction basin. A clean SVHN
rerun is scheduled (see supplementary).}
\end{table}

\subsection{Gain Scaling Law}
\label{sec:scaling}

Collecting results across all nine benchmarks, we find that the mixed-length
ensemble gain scales approximately linearly with the headroom
$(1 - p_\mathrm{individual})$: harder tasks with lower individual accuracy
receive proportionally larger gains.
A least-squares linear fit of gain on headroom gives slope $0.057$,
intercept $0.011$, Pearson $r=0.753$, $p=0.019$, $R^2=0.57$; a power-law
fit $\text{gain} = c \cdot \text{headroom}^{\alpha}$ gives exponent
$\alpha = 0.85$ (log-space $R^2 = 0.72$). Both fits support the intuition
that the gain is governed by the fraction of samples still correctable
by ensembling.
On CIFAR-100, we note that the mixed ensemble at $0.7353$ is slightly
below the best single-length 5-seed ensemble at $B=75$~epochs
($0.7367$); the mixed gain is positive when compared against the average
of the two same-length ensembles (\emph{mean}($A$-5-seed, $B$-5-seed)
$=0.7296$, mixed $-$ mean $=+0.0057$). CIFAR-100 is the one dataset
where the longer schedule $B$ alone beats the mixed ensemble by a narrow
margin; the mixed ensemble still wins over the matched same-length
baseline.

\begin{figure}[h]
\centering
\includegraphics[width=0.75\textwidth]{fig_scaling_law}
\caption{Ensemble gain vs.\ headroom $(1 - \text{individual accuracy})$
across nine benchmarks. Dashed red: through-the-origin fit
$\text{gain} = 0.086 \cdot \text{headroom}$. Dotted green: ordinary
least-squares fit with intercept ($r=0.753$, $p=0.019$, $R^2=0.57$).
A power-law fit gives exponent $\alpha=0.85$.}
\label{fig:scaling}
\end{figure}

\begin{table}[h]
\centering
\caption{Ensemble gain scales inversely with individual accuracy.}
\label{tab:scaling}
\begin{tabular}{lccc}
\toprule
Dataset & Ind.\ accuracy & Mixed gain & Ratio: gain/(1-ind.) \\
\midrule
FACED 9c EEG           & 0.6581 & +0.0367 & 0.107 \\
SEED-V 5c EEG          & 0.3743 & +0.0340 & 0.054 \\
CIFAR-100 CNN          & 0.6886 & +0.0467 & 0.150 \\
CIFAR-100 ResNet-18    & 0.7707 & +0.0397 & 0.173 \\
CIFAR-10 CNN           & 0.9015 & +0.0238 & 0.242 \\
CIFAR-10 ResNet-18     & 0.9452 & +0.0147 & 0.268 \\
Fashion-MNIST CNN      & 0.9313 & +0.0121 & 0.176 \\
SST-2 DistilBERT       & 0.9062 & +0.0032 & 0.034 \\
MNIST MLP              & 0.9849 & +0.0014 & 0.093 \\
\bottomrule
\end{tabular}
\end{table}

\subsection{Comparison to Existing Ensemble Methods (CIFAR-10)}

To verify that our gains are not just from extra training compute, we
compare \ours{} against three standard diversity-boosting ensemble
methods under a matched compute budget on CIFAR-10 with the same small
CNN: (i)~a standard 10-seed same-length deep ensemble, (ii)~Stochastic
Weight Averaging starting at epoch 30 of a 60-epoch run, and
(iii)~a snapshot ensemble with 5 cyclic-LR snapshots from a single run.
All methods were trained with identical hyperparameters except for the
ensemble construction; \ours{} consumes roughly the same total compute
as the 10-seed deep ensemble ($5\times50 + 5\times75 = 625$ vs.\
$10\times50 = 500$ epochs, a $1.25\times$ increase).

\begin{table}[h]
\centering
\caption{Head-to-head comparison with SWA, snapshot ensembles, and standard
10-seed deep ensembling on CIFAR-10 with the same small CNN.
Same total compute budget.}
\label{tab:cifar10_baselines}
\begin{tabular}{lcc}
\toprule
Method & Test Accuracy \\
\midrule
Single model, 50 epochs, mean of 10 seeds & 0.8992 \\
Standard 5-seed deep ensemble (50ep) & 0.9209 (+0.0217) \\
Standard 10-seed deep ensemble (50ep) & 0.9223 (+0.0231) \\
SWA (60ep with SWA starting at 30) & 0.8918 ($-0.0074$) \\
Snapshot ensemble (1 run, 5 cyclic-LR snapshots) & 0.8946 ($-0.0046$) \\
\rowcolor{blond}
\textbf{\ours{} (5$\times$50ep + 5$\times$75ep, 10 models mixed)} & \textbf{\ciftenacc{} (+0.0261)} \\
\bottomrule
\end{tabular}
\end{table}

\begin{figure}[h]
\centering
\includegraphics[width=0.85\textwidth]{cifar10_baselines_bar}
\caption{CIFAR-10 head-to-head ensemble comparison at matched compute.
\ours{} beats every baseline including the 10-seed standard deep
ensemble. SWA and snapshot ensembling both \emph{hurt} relative to the
single-model mean on this small-CNN trainer, likely because they
average over a narrow region of weight space.}
\label{fig:cifar10_bar}
\end{figure}

\subsection{Theoretical Decomposition}

We repeated the experiment on FACED d6 with saved checkpoints and
decomposed the ensemble gain into (i)~within-group and (ii)~cross-group
contributions.

\begin{table}[h]
\centering
\caption{Ensemble gain decomposition on FACED d6. 100ep and 150ep
individual-seed means are identical within 0.0001; the 150ep ensemble is
stronger because 150ep seeds are more diverse.}
\label{tab:decomposition}
\begin{tabular}{lc}
\toprule
Ensemble & Balanced Acc \\
\midrule
Individual 100ep mean (5 seeds) & 0.6581 \\
Individual 150ep mean (5 seeds) & 0.6581 \\
\midrule
5-seed 100ep ensemble & 0.6798 ($+0.0218$) \\
\textbf{5-seed 150ep ensemble} & \textbf{0.6919 ($+0.0338$)} \\
\textbf{10-model mixed} & \textbf{\facedacc{} ($+0.0367$)} \\
\midrule
Mixed gain over 100ep alone & $+0.0150$ \\
\bottomrule
\end{tabular}
\end{table}

\paragraph{Decision-boundary analysis.}
Of the 1{,}932 FACED test samples, 50.9\% have unanimous prediction from
all 10 models; 13.1\% (253 samples) have a \emph{different} group
consensus between 100ep and 150ep. On those 253 samples, the 100ep
consensus accuracy is 28.85\%, the 150ep consensus accuracy is 29.25\%,
and the mixed ensemble accuracy is 39.53\%---a gain of $+10.3$ percentage
points on a subset of 13.1\% of the test set. The predicted overall gain
$10.3 \times 0.131 = 1.35$~pp closely matches the observed $+1.50$~pp.

\subsection{Calibration: NLL, ECE, and Brier}

Ensembles are usually evaluated on calibration as much as
accuracy~\cite{ashukha2020pitfalls}. We
recompute the CIFAR-10 head-to-head on fresh runs, this time saving
full softmax logits, and report accuracy, NLL (negative log-likelihood
in nats), ECE (15-bin expected calibration error), and Brier score.

\begin{table}[h]
\centering
\caption{CIFAR-10 full-calibration comparison (same small CNN, fresh
seeds for this calibration run; reported acc numbers differ slightly
from Table~\ref{tab:cifar10_baselines} due to fresh seeds).
\ours{} improves \emph{all four} metrics over SWA, snapshot, and a
10-seed same-length deep ensemble. Lower is better for NLL, ECE, Brier.}
\label{tab:cifar10_calib}
\begin{tabular}{lcccc}
\toprule
Method & Acc $\uparrow$ & NLL $\downarrow$ & ECE $\downarrow$ & Brier $\downarrow$ \\
\midrule
Single model (50ep) mean of 5 seeds & 0.8904 & 0.3469 & 0.0328 & 0.1637 \\
SWA (60ep)                          & 0.8865 & 0.3537 & 0.0275 & 0.1688 \\
Snapshot (5 cyclic snaps)           & 0.8882 & 0.3372 & 0.0180 & 0.1638 \\
5-seed ens 50ep                     & 0.9058 & 0.2813 & 0.0126 & 0.1389 \\
5-seed ens 75ep                     & 0.9146 & 0.2557 & 0.0104 & 0.1253 \\
5-seed ens 100ep                    & 0.9190 & 0.2479 & 0.0125 & 0.1196 \\
10-seed same-length 50ep            & 0.9078 & 0.2722 & 0.0101 & 0.1355 \\
\rowcolor{blond}
\textbf{\ours{} (5$\times$50ep + 5$\times$65ep)}   & 0.9126 & 0.2607 & 0.0098 & 0.1290 \\
\rowcolor{blond}
\textbf{\ours{} (5$\times$50ep + 5$\times$75ep)}   & 0.9149 & 0.2566 & 0.0107 & 0.1277 \\
\rowcolor{blond}
\textbf{\ours{} (5$\times$50ep + 5$\times$100ep)}  & \textbf{0.9168} & \textbf{0.2489} & \textbf{0.0096} & \textbf{0.1236} \\
\bottomrule
\end{tabular}
\end{table}

\paragraph{Findings.}
(i)~\ours{} with $T_B=100$ achieves the best NLL, the best ECE, and
the best Brier among all methods. (ii)~SWA \emph{hurts} NLL relative
to a single model ($0.3537$ vs.\ $0.3469$) because weight averaging
collapses the ensemble's uncertainty into a point estimate; our method
preserves it by averaging outputs. (iii)~The endpoint-ratio sweep
($T_B \in \{65, 75, 100\}$) shows larger $T_B$ gives modestly better
metrics --- the longer tail phase explores more of the basin.

\subsection{Loss Landscape: Different Basins}

A complementary mechanism test in the style of
Fort~\emph{et al.}~\cite{fort2019deepensembles} is to
check whether the 50-epoch and 75-epoch minima lie in the \emph{same}
loss basin (connected by a monotonic-loss path) or in \emph{different}
basins (separated by a barrier). We train two CIFAR-10 small CNNs with
the \emph{same} random seed (2025) at $T_A=50$ and $T_B=75$ epochs, then
linearly interpolate between their weights
$w(\alpha)=(1-\alpha) w_A + \alpha w_B$ for $\alpha \in [0, 1]$ and
measure the train and test loss at each interpolated point.

\begin{figure}[h]
\centering
\includegraphics[width=0.85\textwidth]{loss_landscape_interp}
\caption{Train/test loss along the linear interpolation between a
$T_A=50$ and a $T_B=75$ CIFAR-10 checkpoint (same seed). Both endpoints
are low-loss minima (train $\approx 0.05$--$0.12$), but the midway
linear path has a \textbf{+0.25 nats train-loss barrier}
and accuracy drops to $0.843$ ($-6$~pp). The two endpoints are not
in the same basin --- they are separated by a loss ridge, directly
confirming that trajectory-length variation drifts weights into a
neighboring local minimum.}
\label{fig:landscape}
\end{figure}

\paragraph{Result.} The train-loss barrier is $+0.2497$ and the
test-loss barrier is $+0.1266$. The test accuracy at $\alpha=0.5$ is
$0.8434$, compared to $0.8905$ at $\alpha=0$ and $0.9018$ at
$\alpha=1$ --- a $-6$~pp drop. This is a textbook ``different basins''
signature. Same seed, different $T_\mathrm{max}$, different minima.

\subsection{Training-Log Analysis: Same Accuracy, Different Trajectories}

To visualize \emph{why} the mixed ensemble works, we plot the full
per-epoch validation curves for all 10 FACED d6 seeds from both groups.
The two trajectories converge to near-identical final accuracy but visit
different regions of the loss landscape at the key epoch 100 crossover
point.

\begin{figure}[h]
\centering
\includegraphics[width=0.75\textwidth]{training_curves_100_vs_150}
\caption{Validation BAcc curves for 100ep and 150ep d6 seeds on FACED. The
100ep trajectory (solid) converges around epoch 100 at BAcc $\approx$0.677.
The 150ep trajectory (dashed) is slightly behind at epoch 100 (still at
BAcc $\approx$0.672, with LR $=0.25\times$ initial), then slowly ascends
to the same final BAcc by epoch 149. The extra 50 epochs add zero mean
accuracy but substantial parameter drift.}
\label{fig:curves}
\end{figure}

\begin{table}[h]
\centering
\caption{Per-epoch val BACC for 100ep and 150ep d6 FACED seeds (mean over
5 seeds). Identical final accuracy, different trajectories.}
\label{tab:trajectories}
\begin{tabular}{lcc}
\toprule
Epoch & 100ep group BAcc & 150ep group BAcc \\
\midrule
25 & 0.6267 $\pm$ 0.0118 & 0.6186 $\pm$ 0.0057 \\
50 & 0.6513 $\pm$ 0.0090 & 0.6516 $\pm$ 0.0046 \\
75 & 0.6718 $\pm$ 0.0046 & 0.6577 $\pm$ 0.0079 \\
99 & 0.6771 $\pm$ 0.0059 & 0.6722 $\pm$ 0.0058 \\
149 & --- & 0.6776 $\pm$ 0.0040 \\
\bottomrule
\end{tabular}
\end{table}

\subsection{Endpoint-Ratio Ablation: Why $T_B=150$?}
\label{sec:endpoint_sweep}

A natural reviewer question is: is $T_B=150$ a tuned sweet spot, or does
the gain hold across a range of endpoint ratios? To answer this directly
we ran the full 5-seed sweep at $T_B \in \{110, 120, 130, 140, 150, 200\}$
on FACED EMOD d6, with $T_A=100$ held fixed in all mixed ensembles.
\cref{tab:endpoint_sweep} and \cref{fig:endpoint_sweep} report, for each
$T_B$, the mean individual BAcc of the $T_B$ group, the $T_B$-only 5-seed
ensemble, and the 10-model mixed ensemble combining $T_A=100$ with each
$T_B$.

\begin{table}[h]
\centering
\caption{FACED EMOD d6 endpoint-ratio ablation. All rows use the
same $T_A=100$ group (mean ind.\ 0.6581, 5-seed ens.\ 0.6798). Mixed
10-model BAcc peaks at $T_B=150$ (our default) and degrades both for
smaller drifts ($T_B \le 140$) and for over-long drifts
($T_B=200$). Source: \texttt{llm\_steering/diag\_p2\_d4\_endpoint\_T*.json}.}
\label{tab:endpoint_sweep}
\begin{tabular}{cccc}
\toprule
$T_B$ & Mean ind.\ $T_B$ & 5-seed ens.\ $T_B$ & \textbf{Mixed 10-model} \\
\midrule
100 (=$T_A$) & 0.6581 & 0.6798 & 0.6798 \\
110 & 0.6612 & 0.6809 & 0.6842 \\
120 & 0.6418 & 0.6740 & 0.6805 \\
130 & 0.6308 & 0.6758 & 0.6778 \\
140 & 0.6292 & 0.6730 & 0.6750 \\
\rowcolor{blond}
\textbf{150} & \textbf{0.6581} & \textbf{0.6919} & \textbf{0.6948} \\
200 & 0.5960 & 0.6425 & 0.6633 \\
\bottomrule
\end{tabular}
\end{table}

\begin{figure}[h]
\centering
\includegraphics[width=0.72\textwidth]{fig_endpoint_sweep}
\caption{Mixed 10-model BAcc on FACED as a function of $T_B$ with
$T_A=100$ fixed. Dashed green: 100ep-only 5-seed ensemble baseline.
The endpoint ratio $\rho = T_B/T_A = 1.5$ is a genuine operating-regime
sweet spot, consistent with the PCA-basin and mode-connectivity
analyses (\cref{sec:basin}).}
\label{fig:endpoint_sweep}
\end{figure}

\paragraph{Interpretation.}
The sweep traces a clear, interpretable curve. For $T_B \in [110, 140]$
the $T_B$ trajectory has not yet spent enough time in the low-LR tail to
reach a distinct basin (mean individual BAcc drops as the cosine
schedule forces a larger tail that the run cannot fully recover, yet the
group has not visited a \emph{separate} minimum either); the mixed ensemble
sits at or slightly above the $T_A$-only baseline. At
$\rho = 1.5$ the mean individual BAcc recovers to $0.6581$
(identical to $T_A$ within $0.0001$), while the run has spent $50$
additional epochs in the tail phase and arrives at a distinct basin
(\cref{sec:basin}); this is where the mixed gain is largest at $+0.015$
BAcc over the $T_A$-only ensemble. At $\rho = 2.0$ the long tail is past
the productive regime and individual accuracy drops to $0.596$, so the
mixed ensemble is pulled below the $T_A$-only baseline. In practice we
recommend $\rho \approx 1.5$; gains at $\rho = 1.1$ are positive but
smaller, and $\rho \ge 2.0$ should be avoided.

\subsection{Loss-Landscape Geometry: The Two Groups Occupy Different Basins}
\label{sec:basin}

The ensemble gain above suggests that the 100ep and 150ep runs end in
\emph{different} local minima. We verify this directly in two
complementary ways: a weight-space PCA of the 10 checkpoints, and a
linear mode-connectivity probe between pairs of minima.

\paragraph{Weight-space PCA.}
Flattening the 10 FACED d6 checkpoints (5 seeds $\times$ 2 endpoints) into
a matrix and taking its first three principal components shows that PC1
captures $>\!99\%$ of the between-checkpoint variance and perfectly
separates the two groups (\cref{fig:pca_basin}). Concretely, the mean
pairwise distance between the two group centroids is $5166$ and the mean
pairwise distance within each group is $1587$, a between/within ratio of
$\mathbf{3.25}$. The 100ep and 150ep runs are thus geometrically
distinct in weight space: they inhabit two separate
neighborhoods, not one broad basin visited at two points along a cosine
schedule.

\begin{figure}[h]
\centering
\includegraphics[width=0.72\textwidth]{fig_pca_basin}
\caption{Weight-space PCA of all 10 FACED d6 checkpoints. PC1 accounts
for $>\!99.9\%$ of variance and fully separates the 100ep and 150ep
groups. Between/within ratio $= 3.25$. Source:
\texttt{reports/p2\_pca\_basin.json}.}
\label{fig:pca_basin}
\end{figure}

\paragraph{Parameter-space averaging fails; prediction-space averaging works.}
To check whether the two groups' minima are connected by a low-loss
linear path (same basin) or separated by a ridge (different basins), we
linearly interpolate between cross-group pairs
$\theta(\lambda) = (1-\lambda)\theta_{100} + \lambda\theta_{150}$ for
$\lambda \in \{0, 0.1, \ldots, 1\}$ and evaluate val and test loss at
each $\lambda$ (\cref{fig:mode_conn}, red curve). The interpolated
weights produce \emph{near-chance} performance for every
$\lambda$ (test loss $\approx 2.18$ nats, test BAcc $0.24$--$0.26$ vs.\
$0.66$ at the endpoints): the linear path between a 100ep and a 150ep
minimum is flat at chance, indicating that the two endpoints are not
even connected by a low-loss straight line in weight space. For reference,
within-group interpolation ($\lambda$-path between two 100ep seeds or
two 150ep seeds, blue/orange dashed) traces the usual mode-connectivity
barrier: test loss peaks near $\lambda=0.5$ at $\approx 2.0$ and
recovers to the endpoint value.

This is a useful negative: it tells us unambiguously that the
\ours{} gain comes from \textbf{prediction-space averaging}, not from
any weight-space convexity between trajectories. Averaging softmax
outputs succeeds precisely where averaging weights fails, because
softmax averaging is agnostic to the geometry of the weight space and
only cares about the per-sample class distributions.

\begin{figure}[h]
\centering
\includegraphics[width=0.95\textwidth]{fig_mode_connectivity}
\caption{Linear interpolation loss on FACED. Solid red:
$\theta_{100}\!\rightarrow\!\theta_{150}$ (cross-group, different seeds);
loss stays at chance level ($\approx 2.18$ nats) along the entire path,
confirming the two groups live in different basins. Dashed blue/orange:
within-group 100ep-to-100ep and 150ep-to-150ep interpolations; these
show the expected $\sim\!1.0$-nat mode-connectivity barrier. Takeaway:
weight averaging between groups is not a meaningful ensemble
substitute; prediction-space averaging preserves the gain.
Source: \texttt{reports/p2\_mode\_connectivity.json}.}
\label{fig:mode_conn}
\end{figure}

\paragraph{Disagreement decomposition (visual).}
\cref{fig:disagree} shows the accuracy of each estimator on the 253
disagreement samples (13.1\% of FACED test). The mixed 10-model
ensemble (rightmost, yellow) achieves $39.53\%$ accuracy on this
subset, $+10.68$~pp above the 100ep consensus and $+10.28$~pp above the
150ep consensus; the gain is cleanly localized to the disagreement set.

\begin{figure}[h]
\centering
\includegraphics[width=0.78\textwidth]{fig_disagreement_decomposition}
\caption{Per-method accuracy on the 253 disagreement samples
(13.1\% of FACED test). The two group consensuses are nearly tied
($28.85\%$ vs.\ $29.25\%$); each group's 5-seed ensemble recovers a few
percentage points; the 10-model mixed \ours{} ensemble recovers the
most. Source: \texttt{reports/p2\_disagreement\_deep\_dive.json}.}
\label{fig:disagree}
\end{figure}
